El siguiente paso natural es fortalecer la conexi\'on entre incompletitud topol\'ogica y completitud observacional cuantitativa. Algunas extensiones prioritarias son:

\begin{itemize}
\item inyecci\'on-recuperaci\'on por m\'etodo de descubrimiento;
\item funciones de completitud espec\'ificas por estrella y por instrumento;
\item simulaciones N-body para evaluar estabilidad din\'amica de intervalos propuestos;
\item integraci\'on expl\'icita de excentricidad e inclinaci\'on;
\item uso de cat\'alogos y series temporales de TESS, Kepler, Gaia y seguimiento RV;
\item validaci\'on con cat\'alogos futuros para verificar si sistemas prioritarios reciben nuevas detecciones o mejores caracterizaciones;
\item comparaci\'on con arquitecturas compactas y sistemas cercanos a resonancia;
\item construcci\'on de rankings observacionales orientados a propuestas de seguimiento.
\end{itemize}

Tambi\'en es deseable generar fichas narrativas por sistema prioritario, donde se integren gap orbital, soporte topol\'ogico, an\'alogos globales, detectabilidad relativa y calidad de datos. Esa salida facilitar\'ia traducir el ranking del pipeline a decisiones observacionales concretas y auditables.
